\begin{recipebox}{Satay Chicken \& Red Lentil Curry}{4}{15 min prep + 6--7 hr (low)}

A creamy peanut and coconut curry in which the red lentils dissolve into the
sauce. Serve over rice. Roughly 40--43\,g protein per serve.

\begin{ingredients}
  \item[] \textbf{Slow cooker base}
  \item \nicefrac{1}{2} cup dry split red lentils, rinsed
  \item 1 onion, diced
  \item 3 garlic cloves, crushed
  \item 1 tsp grated ginger
  \item 400\,ml lite coconut milk
  \item 1 cup chicken stock
  \item 2 tbsp peanut butter
  \item 2 tbsp reduced-salt soy sauce
  \item 1 tbsp curry powder
  \item 500\,g chicken breast or tenderloins, left whole
  \item[] \textbf{Last 20--30 minutes}
  \item 400\,g frozen stir-fry vegetables
  \item[] \textbf{To finish \& serve}
  \item 1 tbsp lime or lemon juice
  \item Chilli flakes (optional)
  \item 1\nicefrac{1}{4} cups dry rice
\end{ingredients}

\begin{method}
  \item \textbf{Build the base.} Put the lentils, onion, garlic, ginger, coconut milk, stock, peanut butter, soy sauce and curry powder into the slow cooker. Mix well so the peanut butter and lentils are spread through the liquid.
  \item \textbf{Add the chicken.} Nestle the chicken into the sauce.
  \item \textbf{Slow cook.} Cook on Low for 6--7 hours (check tenderloins after 5--6). It's fine to drop to Keep Warm afterwards.
  \item \textbf{Shred.} Shred the chicken and stir it back through the sauce.
  \item \textbf{Add the veg.} Stir in the frozen vegetables for the final 20--30 minutes.
  \item \textbf{Finish.} Stir in the lime or lemon juice, and loosen with water if the sauce is thicker than you like.
  \item Cook the rice separately and divide everything into four servings.
\end{method}

\begin{notelist}
  \item \textbf{Advance prep:} dice the onion and defrost the chicken in the fridge the night before.
  \item For a quicker Instant Pot satay, see Thai Peanut Chicken in the Instant Pot chapter.
\end{notelist}

\end{recipebox}
